\scanpage{7}
$\underline{\omega}_A$ est muni d'un élément canonique, que je note $\omega_A$, caractérisé par
\[
\left\{\begin{aligned}
&\omega_A\in\underline{\omega}_A,\\
&\alpha_A(\omega_A)=e_M\qquad\text{ou encore}\qquad\alpha_A=e_M\wedge\omega_A
\end{aligned}\right.\tag*{(14)}
\]
Considérons aussi
\[
\begin{aligned}
\widetilde S_M=M^*&\overset{\mathrm{def}}{=}\left\{x\in M\ \middle|\ \begin{gathered}
x(s)\ne0\text{ dans }M(s)\\ \text{pour tout }s\in\mathrm{Spec}\,\Lambda\end{gathered}\right\}\\
&=\left\{x\in M\mid M/\Lambda x\text{ est un }\Lambda\text{-module inversible}\right\}
\end{aligned}\tag*{(15)}
\]
On a \deleted{une application} \inserted{donc l'inclusion}
\[
\widetilde F_M\subset\unsure{\mathfrak P|\mathscr P}_3(\widetilde S_M=M^*)\tag*{(\deleted{\unsure{6|(}}16)}
\]
et on a en fait
\[
\left\{\begin{aligned}
&\text{Si }\mathrm{Pic}(\Lambda)=0\qquad\text{p.ex. si }\Lambda\ \tfrac12\text{ local}\\
&\text{alors}\quad\widetilde S_M=M^*=\{x\in M\mid\exists\, A\in\widetilde F_M,\ x\in A\}\\
&\text{i.e.}\quad\widetilde S_M=\bigcup_{A\in\widetilde F_M} A
\end{aligned}\right.\tag*{(17)}
\]
Ainsi, $\widetilde S_M$ apparaît comme la réunion d'une ensemble \inserted{$\widetilde F_M$} de parties à trois éléments, et on voit tout de suite que si $\{x,y\}\in\mathfrak P_2(\widetilde S_M)$, alors il existe au plus \uline{un} $A\in\widetilde F_M$ le contenant.
% archive sheet number: 6
